\section{Monte Carlo Design}
\label{chap:monte-carlo-design}
\subsection{Baseline and controlled relevance extension}
\label{sec:monte-carlo-design}
The baseline follows the executable high-dimensional design of \textcite{chen_huang_tien_2021}, with independent observations and 100 controls. The outcome and assignment disturbances satisfy
\begin{equation}
  \begin{pmatrix}
    u_i \\
    \epsilon_i
  \end{pmatrix}
  \sim
  \mathcal{N}\!\left(
    \begin{pmatrix} 0 \\ 0 \end{pmatrix},
    \begin{pmatrix} 1 & 0.3 \\ 0.3 & 1 \end{pmatrix}
  \right).
  \label{eq:dgp-disturbances}
\end{equation}

\begin{equation}
  x_{ji}\overset{\mathrm{ind}}{\sim}\mathcal{N}(0,1),
  \quad j=1,\ldots,100,
  \qquad
  z_{1i},z_{2i},v_{1i},v_{2i}
  \overset{\mathrm{ind}}{\sim}\mathcal{N}(0,1).
  \label{eq:dgp-primitives}
\end{equation}

\begin{equation}
  X_{ji}=\Phi(x_{ji}),
  \qquad j=1,\ldots,100,
  \label{eq:dgp-controls}
\end{equation}
The remaining primitives are mutually independent and independent of $(u_i,\epsilon_i)$. The controls $X_{ji}=\Phi(x_{ji})$ are independent Uniform$(0,1)$ variables. The observed instruments are
\begin{equation}
  Z_{1i}=z_{1i}+v_{1i}+X_{2i}+X_{3i}+X_{4i},
  \label{eq:dgp-instrument1}
\end{equation}

\begin{equation}
  Z_{2i}=z_{2i}+v_{2i}+X_{7i}+X_{8i}+X_{9i}+X_{10i}.
  \label{eq:dgp-instrument2}
\end{equation}
The released code uses transformed $X$ controls in these equations, although the source paper displays corresponding terms as lowercase $x$. We retain the executable construction. The baseline treatment and outcome are
\begin{equation}
  D_i=\Phi(z_{1i}+z_{2i}+\epsilon_i).
  \label{eq:dgp-treatment-baseline}
\end{equation}

\begin{equation}
  Y_i=1+D_i+5\sum_{j=1}^{7}X_{ji}+D_i u_i.
  \label{eq:dgp-outcome}
\end{equation}
The correlation of $u_i$ with $\epsilon_i$ generates endogeneity through the treatment index and the interaction $D_i u_i$. Seven controls enter the outcome directly, and the union of outcome and instrument controls is the ten-control oracle set. At fixed $d\in(0,1)$,
\begin{equation}
  Y_i(d)=1+5\sum_{j=1}^{7}X_{ji}+d(1+u_i).
  \label{eq:dgp-potential-outcome}
\end{equation}

\begin{equation}
  Q_{Y_i(d)\mid X_i}(\tau)
  =1+5\sum_{j=1}^{7}X_{ji}
  +d\left[1+\Phi^{-1}(\tau)\right].
  \label{eq:dgp-structural-quantile}
\end{equation}

\begin{equation}
  \alpha_0(\tau)=1+\Phi^{-1}(\tau).
  \label{eq:dgp-alpha-true}
\end{equation}
The five structural effects are approximately $-0.282$, $0.326$, $1.000$, $1.674$, and $2.282$ in increasing quantile order.

\label{sec:instrument-relevance-extension}
The original design fixes relevance. To reduce the excluded signal without simultaneously changing treatment dispersion, introduce $w_i\sim\mathcal N(0,1)$ independently of all baseline variables and set
\begin{equation}
  d_i^*(\kappa)
  =\kappa(z_{1i}+z_{2i})
  +\sqrt{2(1-\kappa^2)}\,w_i
  +\epsilon_i,
  \label{eq:dgp-kappa-index}
\end{equation}

\begin{equation}
  D_i(\kappa)=\Phi\!\left(d_i^*(\kappa)\right).
  \label{eq:dgp-treatment-kappa}
\end{equation}

\begin{equation}
  \operatorname{Var}\!\left[d_i^*(\kappa)\right]
  =2\kappa^2+2(1-\kappa^2)+1=3.
  \label{eq:dgp-kappa-variance}
\end{equation}

\begin{equation}
  \operatorname{Cov}\!\left(Z_{1i},d_i^*(\kappa)\right)
  =\operatorname{Cov}\!\left(Z_{2i},d_i^*(\kappa)\right)
  =\kappa.
  \label{eq:dgp-kappa-relevance}
\end{equation}

\begin{equation}
  \operatorname{Cov}\!\left(u_i,d_i^*(\kappa)\right)
  =\operatorname{Cov}(u_i,\epsilon_i)=0.3.
  \label{eq:dgp-kappa-endogeneity}
\end{equation}
Joint normality implies $d_i^*(\kappa)\sim\mathcal N(0,3)$ and an invariant joint distribution of $(u_i,d_i^*(\kappa))$, hence of $(u_i,D_i(\kappa))$. The construction preserves the marginal treatment distribution, the relationship generating endogeneity, the outcome equation, structural effects, controls, and instrument construction. It deliberately changes treatment--instrument relationships. Appendix~\ref{app:kappa-derivation} derives the compensating coefficient and these invariances. At $\kappa=1$ the baseline is recovered; the conceptual endpoint $\kappa=0$ removes the primitive excluded signal. The smallest simulated value is positive, so it is not a formal non-identification experiment. The parameter orders designed relevance and is neither an $F$-statistic nor a universal measure of quantile-specific identification.

\subsection{Design dimensions and performance measures}
\label{sec:performance-measures}
\begin{equation}
  \begin{aligned}
    n&\in\{500,1000\},
    &\tau&\in\{0.10,0.25,0.50,0.75,0.90\},\\
    \kappa&\in\{1,0.50,0.25,0.10\},
    &R&=500.
  \end{aligned}
  \label{eq:simulation-design}
\end{equation}
The 40 designs yield 120 estimator--design cells. Within each replication, primitive draws are reused across relevance settings and estimators use the same observed sample. This common-random-number design improves comparison but does not eliminate Monte Carlo uncertainty. Inference uses a 95\% level. A separate 100-replication calibration reports nested OLS $F$-statistics and partial $R^2$ for adding both instruments to a regression of treatment on an intercept and the ten oracle controls; these are descriptive diagnostics without an IVQR threshold.

For replication $r$, point accuracy is summarized by
\begin{equation}
  \begin{aligned}
    \operatorname{Bias}(\tau)
      &=\frac{1}{R}\sum_{r=1}^{R}\left[\alpha_0(\tau)-\widehat{\alpha}_r(\tau)\right], \\
    \operatorname{MAE}(\tau)
      &=\frac{1}{R}\sum_{r=1}^{R}\left|\widehat{\alpha}_r(\tau)-\alpha_0(\tau)\right|, \\
    \operatorname{RMSE}(\tau)
      &=\left\{\frac{1}{R}\sum_{r=1}^{R}\left[\widehat{\alpha}_r(\tau)-\alpha_0(\tau)\right]^2\right\}^{1/2}.
  \end{aligned}
  \label{eq:performance-point}
\end{equation}
MAE and RMSE denote mean absolute error and root mean squared error. The preserved bias convention is truth minus estimate: positive bias indicates underestimation. Coverage and power are evaluated directly at the unrounded true value and specified false candidates, independently of grid membership:
\begin{equation}
  C_r(\tau)
    =\mathbf{1}\!\left\{W_{n,r}\!\left(\alpha_0(\tau)\right)\leq c_{0.95}\right\},
  \qquad
  \widehat{\operatorname{Coverage}}(\tau)
    =\frac{1}{R}\sum_{r=1}^{R}C_r(\tau).
  \label{eq:performance-coverage}
\end{equation}

\begin{equation}
  \widehat{\operatorname{Size}}(\tau)
    =\frac{1}{R}\sum_{r=1}^{R}
      \mathbf{1}\!\left\{W_{n,r}\!\left(\alpha_0(\tau)\right)>c_{0.95}\right\}
    =1-\widehat{\operatorname{Coverage}}(\tau).
  \label{eq:performance-size}
\end{equation}

\begin{equation}
  a_{\Delta}(\tau)=\alpha_0(\tau)+\Delta,
  \qquad
  \Delta\in\{-1.00,-0.50,-0.25,0.25,0.50,1.00\}.
  \label{eq:performance-alternatives}
\end{equation}

\begin{equation}
  \widehat{\operatorname{Power}}(\Delta,\tau)
    =\frac{1}{R}\sum_{r=1}^{R}
      \mathbf{1}\!\left\{W_{n,r}\!\left(a_{\Delta}(\tau)\right)>c_{0.95}\right\}.
  \label{eq:performance-power}
\end{equation}
For DML, direct truth and stored-profile evaluations can use different pivotal draws even at the same candidate. Reported coverage is the repeated-sample acceptance frequency of the direct test, rather than membership reconstructed from a stored grid profile. Raw power uses the common nominal cutoff without size adjustment.

The cutoff and common primary grid are
\begin{equation}
  c_{0.95}=\chi^2_{2,0.95}=5.991464547\ldots.
  \label{eq:wn-critical-value}
\end{equation}

\begin{equation}
  \begin{aligned}
    A_0&=[-1,3],
    &h&=0.1,\\
    G_0&=\{a_j=-1+(j-1)h:j=1,\ldots,41\},\\
       &=\{-1,-0.9,\ldots,2.9,3\}.
  \end{aligned}
  \label{eq:primary-grid}
\end{equation}
The four-unit domain and 0.1 resolution retain the original executable grid for all designs. All five true effects lie strictly inside it. Grid point estimates include discretization, endpoint, and first-tie effects; the grid is a computational domain, not an asserted theoretical parameter space.

Let $I_r(a_j)=\mathbf1\{W_{n,r}(a_j)\le c_{0.95}\}$. The accepted-set measure is
\begin{equation}
  M_r(A_0)
  =\sum_{j=1}^{40}\frac{a_{j+1}-a_j}{2}
    \left[I_r(a_j)+I_r(a_{j+1})\right].
  \label{eq:accepted-measure}
\end{equation}
This trapezoidal sum of binary acceptance indicators ranges from zero to four. Adjacent intervals contribute 0, 0.05, or 0.10; the calculation totals separated accepted components without solving for exact critical-value crossings. It is neither the span between extreme accepted values nor unrestricted confidence-region length. Full-grid acceptance requires acceptance at all 41 points; either-boundary acceptance requires acceptance at at least one endpoint. Neither event establishes acceptance between all points or outside the domain.

Appendix~\ref{app:numerical-domain-diagnostics} examines $A_1=[-3,5]$ and $A_2=[-5,7]$ at spacing 0.1 in a separate 100-replication diagnostic with $n=1000$, all quantiles and estimators, and $\kappa\in\{1,0.25,0.10\}$. It reuses saved samples and interior profiles while adding candidates. Coverage and power remain separate direct evaluations. This bounded diagnostic assesses numerical localization over investigated ranges only.
